%% kapa: exam=1 round=26 year=2015 period=2 form=A subject=appraisal_law subject_ko=부동산관계법규 seq=37 number=37
%% source: Q-Net 감정평가사 기출문제 (cst00309) artlSeq=5138001 file=extracted/1008/5138001/제26회_감정평가사 1차_기출/제26회_감정평가사 1차_기출/2교시_감평_A형.pdf
%% conversion: kordoc 4.24.51 -> tools/gen1.py (auto); review=proofread-sonnet; flags=none
\begin{question}{37}
부동산등기법령상 등기한 권리의 순위에 관한 설명으로 옳지 않은 것은?
\begin{choices}
  \item 같은 부동산에 관하여 등기한 권리의 순위는 법률에 다른 규정이 없으면 등기한 순서에 따른다.
  \item 등기의 순서는 등기기록 중 같은 구에서 한 등기 상호간에는 순위번호에 따른다.
  \item 등기의 순서는 등기기록 중 다른 구에서 한 등기 상호간에는 접수번호에 따른다.
  \item 부기등기의 순위는 원칙적으로 주등기의 순위에 따른다.
  \item 같은 주등기에 관한 부기등기 상호간의 순위는 후 (後) 부기등기가 선(先) 부기등기에 우선한다.
\end{choices}
\end{question}
\answer{5}
